%\begin{center}
%\large \bf \runtitulo
%\end{center}
%\vspace{1cm}
\chapter*{\runtitulo}

\noindent Esta tesis estudia el diseño e implementación de un \zkBridge para \Cardano. El objetivo es autenticar un evento ocurrido en una cadena de origen y habilitar una acción correspondiente en una cadena destino, sin depender de un operador centralizado y respetando los límites reales de ejecución \Onchain. Para ello se propone una arquitectura que separa una capa de actualización, encargada de mantener estado remoto autenticado, y una capa de aplicación, encargada de consumir esa evidencia para autorizar el flujo $\LockTx \rightarrow \MintTx$.

La construcción evita verificar directamente el consenso completo de \Cardano dentro de un circuito. En su lugar, utiliza \Mithril como mecanismo de certificación de estado, una verificación \STM particionada en dos fases sobre Aiken y argumentos \Groth auxiliares para pertenencia de transacciones y actualización de un \SMT anti-replay. El prototipo muestra que esta composición puede expresarse en el modelo \eUTxO mediante \UTxOs canónicos, NFTs de estado, \ReferenceScripts y \ProofReceipt, manteniéndose dentro de los presupuestos de unidades de ejecución disponibles.

El resultado principal es una prueba de viabilidad para una instancia isomorfa de \Cardano: no un bridge productivo heterogéneo completo, sino una arquitectura verificable que delimita con precisión sus supuestos de seguridad, sus invariantes \Onchain y los componentes que quedan como trabajo futuro.

\bigskip

\noindent\textbf{Palabras clave:} \zkBridge, \Cardano, \Mithril, \eUTxO, pruebas de conocimiento cero.
